% day2_seg2_ranproc2/state_space_aug.tex
\begin{frame}
  \onslide<1->{
	\frametitle{State Space Augmentation} 
	\framesubtitle{Motivation}
}
\begin{itemize}
  \onslide<2->{
 \item Kalman filter theory assumes the only remaining forcing function is white process noise
}
  \onslide<3->{
 \item The most frequently used models can all be understood as differential equations (or difference equations) 
 driven by white noise
}
  \onslide<4->{
 \item If assumption of the only forcing function being white is not satisfied due to presence of colored noise, 
 then the colored noise can be whitened
}
  \onslide<5->{
 \item Otherwise, new states need to be added:  this is called state vector or state space \textbf{augmentation}
}
  \onslide<6->{
 \item We'll see an example of state space augmentation later when we discuss filter maintenance
}
\end{itemize}
\end{frame}

\begin{frame}
  \onslide<1->{
	\frametitle{State Space Augmentation} 
	\framesubtitle{General Example:  Augmenting a 2-State Model}
}
\begin{itemize}
  \onslide<2->{
\item Given 2-state model:  
\[
\begin{bmatrix}
\dot{x}_1\\
\dot{x}_2

\end{bmatrix} =
\begin{bmatrix}
a_{11} & a_{12}\\
a_{21} & a_{22}
\end{bmatrix}
\begin{bmatrix}
x_1\\
x_2
\end{bmatrix}+
\begin{bmatrix}
w_1\\
w_2
\end{bmatrix}
\]  
}
  \onslide<3->{
\item  Suppose the "noise" is not actually white;  perhaps there is some long-term drift as exhibited by random walk
}
  \onslide<4->{
\item   This means that $\dot{w}_i=\epsilon_i$, where $\epsilon_i$ is white
}
  \onslide<5->{
\item We augment with new states $x_3=w_1,x_4=w_2$ as follows: 
\[
\begin{bmatrix}
\dot{x}_1\\
\dot{x}_2\\
\dot{x}_3\\
\dot{x}_4

\end{bmatrix} =
\begin{bmatrix}
a_{11} & a_{12} & 1 & 0\\
a_{21} & a_{22} & 0 & 1\\
0 & 0 & 0 & 0\\
0 & 0 & 0 & 0
\end{bmatrix}
\begin{bmatrix}
x_1\\
x_2\\
x_3\\
x_4
\end{bmatrix}+
\begin{bmatrix}
0 & 0 & 0 & 0\\
0 & 0 & 0 & 0\\
0 & 0 & 1 & 0\\
0 & 0 & 0 & 1

\end{bmatrix}
\begin{bmatrix}
0\\
0\\
\epsilon_1\\
\epsilon_2
\end{bmatrix}
\]  
}
\end{itemize}
\end{frame}
\begin{frame}
  \onslide<1->{
	\frametitle{State Space Augmentation} 
	\framesubtitle{}
}
\begin{itemize}
  \onslide<2->{
\item   If original measurement was, say,
\[
\begin{bmatrix}
z_1\\
z_2
\end{bmatrix}=
\begin{bmatrix}
1 & 0\\
0 & 1
\end{bmatrix}
\begin{bmatrix}
x_1\\
x_2
\end{bmatrix}+
\begin{bmatrix}
v_1\\
v_2
\end{bmatrix}
\]
}
  \onslide<3->{
\item  It now becomes, after augmentation: 
\[
\begin{bmatrix}
z_1\\
z_2
\end{bmatrix}=
\begin{bmatrix}
1 & 0 & 0 & 0\\
0 & 1 & 0 & 0
\end{bmatrix}
\begin{bmatrix}
x_1\\
x_2\\
x_3\\
x_4
\end{bmatrix}+
\begin{bmatrix}
v_1\\
v_2
\end{bmatrix}
\]
}
\end{itemize}
\end{frame}
\begin{frame}
  \onslide<1->{
	\frametitle{State Space Augmentation} 
	\framesubtitle{How can you tell when to augment your state vector?}
}
\begin{itemize}
  \onslide<2->{
\item To determine if state vector augmentation is needed, look at the pre-fit residual sequence:  $\nu(k)=\innovation{k}$  
}
  \onslide<3->{
\item Examine its autocorelation function $R_{\nu\nu}$ (or autospectal density)  
}
  \onslide<4->{
\item Determine the correlation time constant and augment the state vector with a Gauss-Markov process with the appropriate
time constant in the feedback gain 
}
  \onslide<5->{
\item Alternatively, you can examine the autospectal density graph and eyeball a correlation constant (See Grewal handouts.)  
}
\end{itemize}
\end{frame}

